% Issue #81, part "Conclusion and future work"; owner: vks-g. Input from reports/phase2/main.tex.
% Numbers repeat Sections results, validation and hybrid; sources are listed there.

\section{Conclusion and Future Work}
\label{sec:conclusion}

Phase~2 compared deep sequence models, machine learning on engineered
features and two hybrids on C-MAPSS FD001, under one leakage-controlled
protocol, with a training cap chosen from the training engines and every
choice fixed by cross-validation before the official test was read once. The
answers to the three questions of Section~\ref{sec:introduction} are as
follows.
\begin{enumerate}
\item \emph{Sequence models against engineered features: a tie.} The LSTM,
temporal CNN and Transformer reach test RMSEs of 15.60 to 15.95, and XGBoost on
the engineered features 15.71. What deep learning adds is that it matches the
engineered features from raw sensor windows. Near failure, the networks are as
accurate as the trees on engineered features (5.1 to 7.6 cycles), while ML on
raw inputs is two to four times worse.
\item \emph{Which encoder: none.} Recurrence, dilated convolution and
attention lie within 0.35 cycles of each other on the test, about their seed
spread. They differ in bias rather than accuracy: the LSTM and Transformer
predict late on average, which the NASA score punishes.
\item \emph{Hybrids: a null result.} The two-stage hybrid has the lowest test
RMSE ($15.18\pm0.55$) and NASA score ($356\pm22$), and gated fusion varied less
across seeds than the LSTM in cross-validation. Neither passes the value test
fixed in advance on the official test, so we do not claim that either adds
value on FD001.
\end{enumerate}
The pipeline mattered more than the architecture. Removing the cap costs every
model 7 to 15 cycles, and raw instead of engineered inputs cost the ML models
1.9 to 4.6. With 100 test engines, models within about a cycle of each other
cannot be ranked, and cross-validation ranks them no better.

Five directions follow from these results.
\begin{itemize}
\item \emph{Seed ensembles.} Averaging five seeds lowered the test RMSE of
Hybrid~B to 14.71, below any single seed. It should be fixed in advance and
evaluated in cross-validation before it is trusted.
\item \emph{Multi-regime subsets.} FD002 and FD004 reused the FD001 settings.
A cap of 150 and longer windows, as in Gupta's best
configurations~\cite{Gupta2026auditing}, chosen on their own training engines,
are the first things to test. The hybrid may also show more value where
operating regimes make features harder to engineer.
\item \emph{More test engines.} The narrowest paired interval spans 1.8
cycles, so gains of the size reported for feature fusion on FD001 cannot be
resolved with 100 engines. Larger benchmarks would give these comparisons the
power they need.
\item \emph{Uncertainty and asymmetric costs.} Maintenance needs a lower bound
on the RUL, not only a point estimate. Predictive intervals~\cite{Sharma2025uncertainty}
and the asymmetric loss, which cut the NASA score by 12\% at unchanged RMSE in
the ablations, both target the cost of late predictions directly.
\item \emph{The cap's blind spot.} Every capped model predicts early for
engines above the cap, which carries most of the error on FD002 and FD004.
A target that keeps some information above the cap would address the largest
remaining source of error.
\end{itemize}
